\scanpage{122}
\margin{77}
or, what comes to the same thing, the values
\[
\frac{\Phi_0}{\Phi_k}(Q),\ldots,\frac{\Phi_T}{\Phi_k}(Q)
\]
from which it follows that
\[
(\Phi_0^{(i)}(Q'),\ldots,\Phi_T^{(i)}(Q'))=(x_0(Q),\ldots,x_T(Q))
\]
\margin{p. 210} and
\[
Q=f_nQ'
\]

It is easy to show that $f_n$ is a continuous mapping (in Zariski's topology) and even a regular one. The latter follows from the fact that, in $Q\in V^{(n)}$ ($x_k\ne0$ at $Q$), any regular function $R$ has the form
\[
R=\frac{F_1\left(\dfrac{x_0}{x_k},\ldots,\dfrac{x_T}{x_k}\right)}{F_2\left(\dfrac{x_0}{x_k},\ldots,\dfrac{x_T}{x_k}\right)}
\]
where $F_1$ and $F_2$ are polynomials and $F_2(Q)\ne0$; the corresponding function on $V$ will be
\margin{\unsure{1)}}
\[
R'=\frac{F_1\left(\dfrac{\Phi_0}{\Phi_k},\ldots,\dfrac{\Phi_T}{\Phi_k}\right)}{F_2\left(\dfrac{\unsure{\phi_0|\Phi_0}}{\unsure{\phi_k|\Phi_k}},\ldots,\dfrac{\unsure{\phi_T|\Phi_T}}{\unsure{\phi_k|\Phi_k}}\right)}
\]
\mspage{121}
